\section{Anomalous boundaries and independent multiplicities}
\label{sec:intrinsic-anomaly}

The CZX construction provides a different comparison: an anomalous
core survives even though the added weights separate exactly.

\subsection{Raw blocks and integer fusion}

Use a two-qubit cell with parity embeddings
$I_f\ket t=\ket{t,t\oplus f}$. Binary addition is denoted by
$\oplus$. The raw vertical blocks act on $\C^2$:
\begin{align}
	B_f^{00} & =\Id,                                 \\
	B_f^{t+1,u+1}
	         & =(-1)^{(t+f)u+t(1-f)}\ket t\bra{1-t},
	\qquad t,u\in\{0,1\}.
	\label{eq:CZX-B}
\end{align}
Cross-block letters vanish. The physical unweighted cell is
$\bigoplus_fB_f/4$, not $\bigoplus_fB_f$.

The transfer map is
\begin{align}
	\mc E_{B_f}
	\begin{pmatrix}a&b\\c&d\end{pmatrix}
	 & =\begin{pmatrix}a+2d&b\\c&d+2a\end{pmatrix}.
\end{align}
Its eigenvalues on $\Id,Z,E_{01},E_{10}$ are $3,-1,1,1$.
Thus $\widehat B_f=B_f/\sqrt3$ are normal. Products of two
letters span all matrix units. The identity letter fixes any
equivalence phase to one, while the $(1,1)$ and $(1,2)$ letters
would require the same $E_{01}$ both to retain and to reverse
sign. Hence the two sectors are inequivalent.

For binary $f,g,c$, define
\begin{align}
	h(f,g,c) & =(f+1+g)(c+g)\pmod2,                        \\
	J_{fg}^c\ket t
	         & =(-1)^{h(f,g,c)t}\ket{t,t\oplus c\oplus g}.
\end{align}
The two ranges partition $\C^2\otimes\C^2$. Substitution gives
\begin{align}
	B_fB_g & =\sum_cJ_{fg}^cB_c(J_{fg}^c)^\dagger.
	\label{eq:CZX-fusion}
\end{align}
For the nonidentity letters, put $v=t\oplus c\oplus g$.
The exponent difference between the concatenated coefficient and
the target coefficient is
\begin{align}
	 & (t+f)v+t(1-f)+(v+g)u+v(1-g)
	-(t+c)u-t(1-c)\nonumber        \\
	 & \qquad=(f+1+g)(c+g)\pmod2,
\end{align}
which is the conjugation sign of $J_{fg}^c$.

The raw algebra has $ff'=e+o$ for every pair. With
$p=(e+o)/4$ and $x=e-o$,
\begin{align}
	p^2 & =p, & px & =xp=0, & x^2 & =0.
\end{align}
There is no identity: $ue=(u_e+u_o)(e+o)$ cannot equal $e$.
On the independent formal symbols $e,o$, this algebra is two-dimensional.
Its closed-operator representation need not be faithful at short lengths:
$O_1(B_0)=O_1(B_1)$. The formal algebra has a one-dimensional nilpotent
ideal, and is not a fusion ring with a tensor unit. Horizontal
positivity is compatible with this because the vertical span is
not assumed to be closed under the physical adjoint.

\subsection{Exact shifted-cell factorization}

Define the complete Bell embeddings
\begin{align}
	\Phi_b\ket j
	        & =2^{-1/2}\sum_{r=0}^1(-1)^{jr}\ket{r,r\oplus b},   \\
	\Lambda & =\lambda_0+\lambda_1,                            &
	w       & =\frac1{8\Lambda},
	\qquad \lambda_0,\lambda_1>0.
\end{align}
After regrouping the complete register into one site, the tensor is
\begin{align}
	C^{\alpha\beta}
	 & =w\sum_{b,f,\gamma}\lambda_b
	\Phi_bB_b^{\alpha\gamma}\Phi_b^\dagger
	\otimes I_fB_f^{\gamma\beta}I_f^\dagger.
\end{align}
The embeddings $V_{bf}^c=(\Phi_b\otimes I_f)J_{bf}^c$ have
eight two-dimensional ranges exhausting the physical site.
Therefore
\begin{align}
	C^{\alpha\beta}
	          & =\sum_{b,f,c}w\lambda_b
	V_{bf}^cB_c^{\alpha\beta}(V_{bf}^c)^\dagger,                  \\
	\eta_c
	          & =w\diag(\lambda_0,\lambda_0,\lambda_1,\lambda_1), \\
	\tr\eta_c & =2w\Lambda=\frac14.
	\label{eq:CZX-eta}
\end{align}
Unlike \cref{eq:category-spectra}, these multiplicities are
identical for both output sectors.

\begin{proposition}[CZX spectator separation]
	\label{prop:CZX-separation}
	Let
	\begin{align}
		\sigma_\lambda
		 & =\frac1{2\Lambda}
		\diag(\lambda_0,\lambda_0,\lambda_1,\lambda_1).
	\end{align}
	There is a fixed on-site unitary $W$, following the regrouping,
	such that
	\begin{align}
		WC^{\alpha\beta}W^\dagger
		 & =\sigma_\lambda\otimes
		\left(\frac14\sum_c I_cB_c^{\alpha\beta}I_c^\dagger\right), \\
		W^{\otimes N}\rho_N(C)(W^\dagger)^{\otimes N}
		 & =\sigma_\lambda^{\otimes N}
		\otimes\rho_{\mathrm{CZX}}^{(2N)}
		\qquad(N\geq1).
		\label{eq:CZX-separation}
	\end{align}
\end{proposition}

\begin{proof}
	The sector coordinates must be converted to physical qubits by the
	permutation $P\ket{c,t}=I_c\ket t=\ket{t,t\oplus c}$.
	Define $W V_{bf}^c\ket t=\ket{b,f}\otimes I_c\ket t$.
	Completeness of the ranges and unitarity of $P$ make $W$ unitary.
	Each coefficient in
	\cref{eq:CZX-eta} becomes
	$w\lambda_b=(\lambda_b/(2\Lambda))(1/4)$, proving the letter
	identity. Horizontal contraction proves the all-length statement.
\end{proof}

The normalized BNT data are
\begin{align}
	\mu_c      & =\sqrt3\,\eta_c,   &
	m_c        & =\frac{\sqrt3}{4},   \\
	\chi_{fgc} & =\frac1{\sqrt3},   &
	c_{fgc}(L) & =3^{-L/2}.
\end{align}
The raw and normalized mass equations are, respectively,
\begin{align}
	\sum_{f,g}\frac14\frac14        & =\frac14,          \\
	\sum_{f,g}\frac1{\sqrt3}
	\left(\frac{\sqrt3}{4}\right)^2 & =\frac{\sqrt3}{4}.
\end{align}
The embeddings are complete, so \cref{lem:channels} gives full-space
RFP channels with no complement branch.

Horizontally, the identity block is scalar. The spin-flip block has
dimension two; its blocked letters are nonzero multiples of
$\ket t(1,(-1)^u)$ and span all two-by-two matrices. Its
physical-trace matrix vanishes. The two blocks have disjoint diagonal
and off-diagonal physical supports, and positive register weights
preserve their spans. The shifted tensor is therefore canonical and
nonsimple for every positive weight pair.

This factorization uses a regrouped cell and changes the canonical
multiplicities. It is therefore stronger than testing positive rescalings
at the original register-crossing cut, and is not contradicted by
length-dependent coefficients in that presentation. No crossing-cut
normalization is needed for the shifted tensor or its
all-positive-weight factorization.

\subsection{Anomaly and bulk interpretation}

For an even ring of $m$ qubits,
\begin{align}
	U_m & =\left(\prod_{j=1}^mCZ_{j,j+1}\right)X^{\otimes m}, \\
	\rho_{\mathrm{CZX}}^{(m)}
	    & =2^{-m}(\Id+U_m).
\end{align}
Bit complementation changes $\sum_jt_jt_{j+1}$ by
$m-2\sum_jt_j$. Hence
\begin{align}
	U_m^2       & =(-1)^m\Id, &
	U_m^\dagger & =(-1)^mU_m.
\end{align}
At even $m$, the operator is a Hermitian involution and is traceless,
so the displayed density operator is positive and trace one.
The two-qubit cell is necessary for an all-length MPDO statement.

On an interval of $m\geq2$ qubits, retaining only its $m-1$ internal
controlled-$Z$ gates gives
\begin{align}
	U_I^2                  & =(-1)^{m-1}Z_1Z_m, \\
	U_IZ_1U_I^{-1}         & =-Z_1,             \\
	Z_1U_IZ_1^{-1}U_I^{-1} & =-\Id.
\end{align}
This endpoint sign expresses the nontrivial symmetry obstruction of
the CZX boundary~\cite{CZX}. The separation in
\cref{prop:CZX-separation} does not remove the anomalous core.
Nor does it prove that the separating circuit preserves an
independently specified symmetry action. That stronger assertion
requires compatibility of the gates with that action.

Encodings of unweighted CZX or CZY boundaries into double-semion
descriptions do not automatically construct a bulk tensor for every
weighted presentation. A positive filter commuting with a boundary
projector is not generally trace preserving on arbitrary inputs.
Neither the fixed-cut coefficients nor the spectator state establish
a new two-dimensional phase or an equivalence in a prescribed
symmetry-preserving class.
